\section{Introduction}
Time series (TS) data are ubiquitous across a wide range of domains, including economics, urban computing, and epidemiology~\cite{sezer2020financial,tabassum2021actionable,kamarthi2021doubt}. Analytical tasks on such datasets hence find broad applications in various real-world scenarios such as energy forecasting, traffic planning, and epidemic policy formulation. Human experts typically complete such Time-Series Analysis (TSA) tasks by integrating multiple modalities of time-series data. For instance, epidemiologists combine numerical data on influenza infections with textual domain knowledge, policies, and reports to predict future epidemiological trends. However, most existing TSA models~~\cite{vaswani2017attention,liu2023itransformer,zhang2022crossformer,liu2022non,zhou2022film,kitaev2020reformer,zhou2021informer,nie2022time} are unimodal, solely using numerical series. 

Recently, with the development of Large Language Models (LLMs), the field of TSA is also undergoing an exciting transformative moment with the integration of natural language~~\cite{yu2023temporal, li2024frozen}. Existing LLM-based TSA methods incorporate endogenous text derived from numerical series, such as linguistic descriptions of statistical information, which has demonstrated promising benefits~~\cite{gruver2024large, jin2023time, cao2023tempo, liu2024lstprompt}. However, the potential of exogenous or auxiliary textual signals—such as information on concurrent events and policies that provide additional context to time series—remains untapped. This observation prompts a crucial question for multimodal TSA: \textbf{Can multimodal TSA models utilize these exogenous textual signals effectively, thereby enhancing current TSA tasks and enabling new applications?}

The primary obstacle in addressing this question lies in the absence of a comprehensive, high-quality multimodal TS dataset, as evidenced by three significant gaps: (1) \textbf{Narrow data domains.} Data characteristics and patterns vary between different domains, such as the periodicity of numerical data and the sparsity of textual data. However,  current multimodal TS datasets~~\cite{dong2024fnspid,finance2,finance3,finance4,cao2023tempo} focus solely on stock prediction tasks in the financial domain, which are unable to represent the diverse data domains. (2) \textbf{Coarse-grained modality alignment.} Existing multimodal TS datasets only ensure that the text and numerical data come from the same domain, such as general stock news and the prices of one specific stock. Clearly, an abundance of irrelevant text diminishes the effectiveness of multimodal TSA. (3) \textbf{Inherent data contamination.} Existing multimodal TS datasets overlook two main reasons of data contamination: (1) Textual data often contains predictions. For example, influenza outlook is a regular section in influenza reports. (2) Outdated test set, particularly the textual data, may have been exposed to LLMs, which are pretrained on vast corpuses. For example, the knowledge cutoff for Llama3-70B is December 2023, which is later than the cutoff dates for most existing multimodal TS datasets. These reasons lead to biased evaluations of general or LLM-based TSA models.

To address the identified gaps, this work aims to introduce a comprehensive, high-quality multimodal TS dataset that spans diverse domains and can be validated through its effectiveness and benefits for TSA. The main contributions of our work are:
\begin{itemize}
    \item \textbf{Pioneering Multi-Domain Multimodal Time-Series Dataset.} We introduce \model, the first multi-domain multimodal time-series dataset that addresses the aforementioned gaps: (1) encompasses 9 primary data domains. (2) ensures fine-grained modality alignment through meticulously selected data sources and rigorous filtering steps. (3) disentangles facts and predictions from text; ensures all cutoff dates are up to May 2024. To the best of our knowledge, \model stands as the inaugural high-quality and comprehensive multimodal time-series dataset. We envision \model offering exciting opportunities to significantly advance time series analysis through multimodal extensions.
    \item \textbf{Pilot Multimodal Time-Series Forecasting Study.} We develop the first-cut multimodal time-series forecasting (TSF) library, \mmtsf, piloting multimodal TSA research based on \model. Our library \mmtsf features an end-to-end pipeline with a seamless interface that allows the integration of any open-source language models with arbitrary TSF models, thereby enabling multimodal TSF tasks. \mmtsf facilitates easy exploration of \model and supports future advancements in multimodal TSA.
    \item \textbf{Extensive Evaluations with Significant Improvement.} We conducted over 1,000 experiments of multimodal TSF on \model using \mmtsf. The multimodal versions outperformed corresponding unimodal versions in 95\% of cases, reducing the mean squared error by an average of over 20\% and up to 40\% in some domains with rich textual data. This significant and consistent improvement demonstrates the high quality of \model, the effectiveness of \mmtsf, and the superiority of multimodal extensions for TSF.
\end{itemize}
We include additional related works in Appendix~\ref{app:related} and limitations in Appendix~\ref{app:limitaion}.
